% !PRINTED 336
\weberpage{336}
%<<CONT
and higher powers of $h$ occur, has a small value. Under suitable
hypotheses, then,
\weq{4}{ \alpha' = \alpha - \frac{f(\alpha)}{f'(\alpha)} }
is to be regarded as a better approximation to the true value of the
root.

If one then replaces $\alpha$ by $\alpha'$, one will obtain in
\[ \alpha'' = \alpha' - \frac{f(\alpha')}{f'(\alpha')} \]
a still better approximation, and so on.

But there still remain here the following two questions to answer:
\begin{enumerate}
\item[1.] \gesperrt{Under what hypotheses is $\alpha'$
  really a better value than $\alpha$?}
\item[2.] \gesperrt{How can one estimate the degree of accuracy that
  one attains in this way?}
\end{enumerate}
These questions were answered by \gesperrt{Fourier}; but in doing so he
makes the following assumption:

\gesperrt{A root of $f(x)$ is enclosed in an interval $(\alpha,\beta)$
that contains no second root.}

\gesperrt{In the interval $(\alpha,\beta)$, $f'(x)$ is different from zero and
$f''(x)$ is also different from zero.}

As regards the latter assumption, that $f''(x)$ is different from zero
in the interval, it is made only in order to obtain conditions for the
application of the method that can be expressed more simply. In itself
the usefulness of the method, given sufficient narrowing of the
interval, does not depend on it. But if $f(x)$ and $f''(x)$ have no
common divisor, and so do not vanish simultaneously, then by narrowing
the interval one can always satisfy the \gesperrt{Fourier} assumption;
and if $f(x)$ and $f''(x)$ do have a common divisor, one can first
separate it off and then apply the approximation methods to the
individual factors of $f(x)$.

One surveys the situation most simply in geometry, by interpreting
$y=f(x)$ as the equation of a plane curve in a rectangular coordinate
system.

The equation
\weq{3}{ y = f(\alpha) + (x-\alpha)f'(\alpha) } % [sic] equation number (3) repeats the earlier (3) used on the previous printed page for f(alpha)+hf'(alpha)=0; reproduced as printed, Weber restarts the numbering for this geometric sub-passage
is the equation of the tangent to the curve at the point $\alpha, f(\alpha)$
%CONT>>
